\section[Continuous mapping]{Continuous mapping \marginpar{\red{Lecture 3}}}
\label{sec: continuous mapping}

An incredibly useful tool in analysis was the concept of continuous functions
that allows the transfer of limits through functions. In this section
we study how the stochastic convergence concepts interact with continuous
functions. In essence we would like that \(f(X_n)\) converges to \(f(X)\)
whenever \(X_n\) converges to \(X\) and \(f\) is continuous.
Our next result shows that this is mostly the case.

\begin{mdframed}[innertopmargin=0pt]
\begin{theorem}[Continuous mapping]
    \label{thm: continuous mapping}
    Let \((\domain, \norm{\cdot})\) and \((\range, \norm{\cdot})\) be normed
    vectorspaces.  Let \((X_n)_{n\in \nat}\) and \(X\) be random variables in
    \(\domain\). Let \(f\colon \domain \to \range\) be almost
    surely continuous in \(X\), that is \(\prob{f \text{ is continuous at } X} =
    1\). Then
    \begin{enumerate}[label=(\roman*)]
        \item\label{it: cont. mapping as} \(X_n \overset{\as}\to X\) implies \(f(X_n) \overset{\as}\to f(X)\).
        \item\label{it: cont. mapping p} \(X_n \overset{p}\to X\) implies \(f(X_n) \overset{p}\to f(X)\).
        \item\label{it: cont. mapping d} \(X_n \overset{d}\to X\) implies \(f(X_n) \overset{d}\to f(X)\).
        \item\label{it: cont. mapping Lp} If \(f\) is \underline{Lipschitz} continuous and \(X_n \overset{L^p}\to X\),
        then \(f(X_n) \overset{L^p}\to f(X)\).
    \end{enumerate}
\end{theorem}
\end{mdframed}

\begin{remark}[Generalization]
    The normed vectorspace is only necessary for \(L^p\) convergence. The other
    implications also hold in metric spaces.
\end{remark}

This continuous mapping result can be very misleading, as one may expect that
\(X_n \overset{d} \to X\) and \(Y_n \overset{d} \to Y\) implies that \(X_n + Y_n
\overset{d} \to X + Y\) since the sum is continuous. This is however
false in general (Exercise \ref{ex: counterexample joint convergence}).
The trap is that the function \((x,y) \mapsto x+y\) may be continuous
but the vector \((X_n, Y_n)\) does not necessarily converge to \((X,Y)\)
if its components do. The second important result of this section
is therefore the following theorem that shows that for most convergence
concepts the convergence of the components is sufficient.

\begin{mdframed}[innertopmargin=0pt]
\begin{theorem}[Joint convergence]
    \label{thm: joint convergence} Let \((X_n)_{n\in \nat}\) and \(X\) be
    random variables in \(\real^\dims\) and  \((Y_n)_{n\in
    \nat}\) and \(Y\) be random variables in \(\real^{\dims'}\). Then
    \begin{enumerate}[label=(\roman*)]
        \item \label{it: joint conv as}
        if \(X_n \overset{\as}\to X\) and \(Y_n \overset{\as}\to Y\),
        then \((X_n, Y_n) \overset{\as}\to (X,Y)\),

        \item \label{it: joint conv Lp}
        if \(X_n \overset{L^p}\to X\) and \(Y_n \overset{L^p}\to Y\),
        then \((X_n, Y_n) \overset{L^p}\to (X,Y)\),

        \item \label{it: joint conv in prob}
        if \(X_n \overset{p}\to X\) and \(Y_n \overset{p}\to Y\),
        then \((X_n, Y_n) \overset{p}\to (X,Y)\),
        \item\label{it: joint conv law slutsky}
        (\textbf{Slutsky})
        if \(X_n \overset{d}\to X\) and \(Y_n \overset{d}\to c\) \underline{for \(c\in \real^{\dims'}\)},
        then \((X_n, Y_n) \overset{d}\to (X,c)\),

        \item\label{it: joint conv law cramer wold}
        (\textbf{Cramér-Wold})
        if \(\scp{t, X_n} \overset{d}\to \scp{t,X}\) for all \(t\in \real^\dims\),
        then \(X_n \overset{d}\to X\).
    \end{enumerate}
\end{theorem}
\end{mdframed}

The statements for convergence in distribution \ref{it: joint conv law slutsky}
and \ref{it: joint conv law cramer wold} are weaker on purpose as the general
statement is false (Exercise \ref{ex: counterexample
joint convergence}).

\begin{remark}[Generalization]
    Except for Cramér-Wold \ref{it: joint conv law cramer wold} these statements
    do not necessarily need to be restricted to \(\real^\dims\). However for
    convergence of the components to imply convergence of the tuple it is
    necessary that the metric on the product space is ``compatible'' with the
    metrics on the components. This is already necessary in the deterministic
    setting (see Exercise~\ref{ex: incompatible metric}). To avoid these
    technicalities we restrict ourselves to this important special case. The
    property we use is
    \[
        \|(x, 0)\| = \|x\|
        \quad\text{and}\quad
        \|(0,y)\| = \|y\|,
    \]
    which is true for \(\real^{\dims + \dims'}\) with the standard Euclidean norms.
    \fxnote{
        Our proofs would probably work directly in Hilbertspaces. Exericise?
    }
\end{remark}

In the rest of this section we prove the Theorems \ref{thm: continuous mapping}
and \ref{thm: joint convergence}.

\begin{proof}[Proof of Theorem \ref{thm: continuous mapping}]
    Let \(D_f\) be the set of discontinuity points of \(f\).
    Then
    \(\prob{X \in D_f} = 0\) by assumption. We have \ref{it: cont. mapping as}
    by
    \[
        \set{X_n \to X} \subseteq \set{f(X_n) \to f(X)} \cup \set{X \in D_f}
    \]
    and consequently
    \[
        \prob[\Big]{f(X_n) \to f(X)}
        \ge \prob[\big]{X_n \to X} - \prob{X \in D_f} = 1.
    \]
    For \ref{it: cont. mapping d} we use Portemanteau's theorem.
    Specifically, for \(f(X_n) \overset{d}\to f(X_n)\) we need
    \begin{equation}
        \label{eq: cont. mapping d}    
        \E{g(f(X_n))} \to \E{g(f(X))}
    \end{equation}
    for all bounded and continuous functions \(g\) \ref{it: b.cont}. To prove
    \eqref{eq: cont. mapping d} we use \ref{it: b.a.s. cont} on \(X_n
    \overset{d}\to X\) and the function \(h=g\circ f\). Observe that \(h\) is
    bounded (as \(g\) is bounded) and almost surely continuous in \(X\) as \(g\)
    is continuous and \(f\) is almost surely continuous in \(X\), thus \ref{it:
    b.a.s. cont} is applicable.
    
    \noindent For \ref{it: cont. mapping Lp} we simply apply Lipschitz continuity of \(f\).
    Let \(K\) be the Lipschitz constant of \(f\), then
    \[\begin{aligned}
        \norm{f(X_n) - f(X)}_{L^p}^p
        &= \E{\norm{f(X_n) - f(X)}^p}
        \\
        \overset{\text{Lip.}}&\le \E{(K \norm{X_n - X})^p}
        = K^p \norm{X_n - X}_{L^p}^p \to 0.
    \end{aligned}
    \]
    Finally, \ref{it: cont. mapping p} is the most involved. For a fixed \(\epsilon\) we define
    \[
        A_\delta \coloneq \set{x \in \domain: \underbrace{x \notin D_f}_{\text{cont. point}}, \underbrace{\exists y \in \domain : \|x-y\|\le \delta, \|f(x) - f(y)\| > \epsilon}_{\delta\text{ not small enough for }\epsilon}}
    \]
    Observe that since \(A_\delta\) only includes continuity points of \(f\) we have \(A_\delta \downarrow \emptyset\) for \(\delta \to 0\).
    By definition of \(A_\delta\) we have the following inclusion of events
    \[
        \set[\big]{\norm{f(X_n) - f(X)} > \epsilon}
        \subseteq \set[\big]{\norm{X_n - X} > \delta} \cup \set{X \in A_\delta} \cup \set{X\in D_f}.
    \]
    Taking the limit superior in \(n\) we thus get
    \[\begin{aligned}
        &\limsup_{n\to \infty} \prob[\big]{\norm{f(X_n) - f(X)} > \epsilon}
        \\
        &\le \underbrace{\lim_{n\to \infty} \prob[\big]{\norm{X_n - X} > \delta}}_{=0}
        + \prob{X \in A_\delta} + \underbrace{\prob{X \in D_f}}_{=0}.
    \end{aligned}
    \]
    Since this is true for all \(\delta > 0\) we have
    \[
        \limsup_{n\to \infty} \prob[\big]{\norm{f(X_n) - f(X)} > \epsilon}
        \le \lim_{\delta \to 0} \prob{X \in A_\delta}
        = \prob[\Big]{ X \in \bigcap_{\delta >0} A_\delta }
        = 0,
    \]
    where we used that \(A_\delta \downarrow \emptyset\). This proves the claim.
\end{proof}


\begin{proof}[Proof of Theorem \ref{thm: joint convergence}]
    For \ref{it: joint conv as} we simply observe
    \[
        \prob[\Big]{\lim_{n\to \infty} (X_n, Y_n) \neq (X,Y)}
        \le \prob[\Big]{\lim_{n\to \infty} X_n \neq X} + \prob[\Big]{\lim_{n\to \infty} Y_n \neq Y} = 0.
    \]
    For \ref{it: joint conv Lp} we use Minkowski's inequality
    \[\begin{aligned}
        \norm{(X_n, Y_n) - (X,Y)}_{L^p}
        &\le \norm{(X_n, Y_n) - (X,Y_n)}_{L^p} + \norm{(X,Y_n) - (X,Y)}_{L^p}
        \\
        &= \norm{X_n - X}_{L^p} + \norm{Y_n - Y}_{L^p} \to 0,
    \end{aligned}\]
    where the last equality uses \(\norm{(x, 0)} = \norm{x}\) in \(\real^{\dims\times \dims'}\) to get
    \[
        \|(X,0)\|_{L^p}^p = \E{\|(X,0)\|^p} \overset{\text{in }\real^{\dims\times \dims'}}= \E{\|X\|^p} = \|X\|_{L^p}^p.
    \]
    For \ref{it: joint conv in prob} we use the stochastic triangle inequality (Exercise \ref{ex: stochastic triangle})
    and again \(\norm{(x, 0)} = \norm{x}\) in \(\real^{\dims\times \dims'}\)
    to obtain
    \[\begin{aligned}
        &\prob[\Big]{\norm{(X_n, Y_n) - (X,Y)} > \epsilon}
        \\
        &\le \prob[\Big]{\norm{(X_n, Y_n) - (X,Y_n)} > \tfrac\epsilon2}
        + \prob[\Big]{\norm{(X,Y_n) - (X,Y)} > \tfrac\epsilon2}
        \\
        &= \prob[\Big]{\norm{X_n - X} > \tfrac\epsilon2}
        + \prob[\Big]{\norm{Y_n - Y} > \tfrac\epsilon2} \to 0.
    \end{aligned}\]
    For \ref{it: joint conv law slutsky} we use that \(Y_n \overset{d}\to c\)
    for some constant \(c\in \real^{\dims'}\) implies \(Y_n \overset{p}\to c\)
    by Theorem \ref{thm: convergence implications}~\ref{it: d to p}. To prove
    \((X_n, Y_n) \overset{d}\to (X,c)\) we show \ref{it: b.cont} from the Portemanteau
    theorem. Thus let \(f\colon \real^{\dims + \dims'} \to \real\) be bounded
    and continuous function. Then \(g(x) \coloneq f(x,c)\) is a bounded
    and continuous function. Since \(X_n \overset{d}\to X\) we therefore have
    \[
        \E{f(X_n, c)} = \E{g(X_n)} \to \E{g(X)} = \E{f(X,c)}.
    \]
    This implies \((X_n, c) \overset{d}\to (X,c)\). With 
    \(\|(X_n, Y_n) - (X_n, c)\| = \|Y_n - c\| \overset{p}\to 0\)
    due to \(Y_n \overset{p}\to c\), we can apply the
    pull-along lemma (Lemma \ref{lem: pull along}) to conclude that \((X_n, Y_n)
    \overset{d}\to (X,c)\).

    We will delay the proof of \ref{it: joint conv law cramer wold} to Section \ref{sec: moment generating functions}
    (Corollary \ref{cor: cramer wold}).
\end{proof}

\begin{lemma}[Pull along]
    \label{lem: pull along}
    \(X_n \overset{d}\to X\) and \(\norm{Y_n - X_n} \overset{p}\to 0\) implies \(Y_n \overset{d}\to X\).
\end{lemma}
\begin{proof}
    We prove \ref{it: b.lip} from the Portemanteau theorem.
    Let \(f\colon \domain \to \real\) be bounded and Lipschitz continuous, then
    \[
        \abs[\big]{\E{f(Y_n)} - \E{f(X)}}
        \overset{\Delta}\le \abs[\big]{\E{f(Y_n)} - \E{f(X_n)}} + \abs[\big]{\E{f(X_n)} - \E{f(X)}}
    \]
    converges to zero as the first term converges to zero by Lemma \ref{lem:
    moving target p to d} and the second term converges to zero by \(X_n
    \overset{d}\to X\).
\end{proof}


\subsection*{Exercises}

\begin{exercise}[Counterexample joint convergence \Coffeecup]
    \label{ex: counterexample joint convergence}
    Let \(X\) be a random variable with \(\prob{X=-1} = \prob{X=1} = \frac12\)
    and define \(X_n \coloneq X\) and \(Y_n \coloneq -X\). Show that \(X_n
    \overset{d}\to X\) and
    \(Y_n \overset{d}\to X\) but 
    \[
        X_n + Y_n \overset{d}{\not\to} X + X = 2X.
    \]
    Conclude that \((X_n, Y_n)\) does not converge to \((X,X)\) in distribution.
    \begin{remark}
        Observe that \(X_n \overset{p}\to X\). It is therefore not sufficient
        that one of the sequences converges in probability as Slutzky (Theorem
        \ref{thm: joint convergence}~\ref{it: joint conv law slutsky}) may
        suggest.
    \end{remark}
    \begin{solution}
        Clearly, \(X_n + Y_n = 0\) for all \(n\). In particular we have
        \[
            \prob{X_n + Y_n \in A} = \ind{A}(0)
            \neq \frac12\ind{A}(2) + \frac12 \ind{A}(-2)
            = \prob{2X \in A}.
        \]
        for \(A= \{0\}\) for example. Since \(\prob{2X = 0} = 0\) this is
        however a requirement for \(X_n + Y_n \overset{d}\to 2X\) by \ref{it:
        s.zero boundary}. Thus \(X_n + Y_n \overset{d}{\not\to} 2X\).  This also
        implies that \((X_n, Y_n)\) does not converge to \((X,X)\) in
        distribution as otherwise we would have \(X_n + Y_n \overset{d}\to X +
        X\) by the continuous mapping theorem.
    \end{solution}
\end{exercise}

\begin{exercise}[Hölder continuos functions \Coffeecup]
    Let \(f\) be a Hölder continuous function. That is
    \[
        \norm{f(x) - f(y)} \le L \norm{x-y}^\alpha
    \]
    Show that \(X_n \overset{L^p}\to X\) implies \(f(X_n) \overset{L^q}\to f(X)\) for \(q = \frac{p}{\alpha}\).
    \begin{solution}
        We have by Hölder continuity of \(f\) that
        \begin{align*}
            \norm{f(X_n)- f(X)}_{L^q}^q
            &= \E{\norm{f(X_n) - f(X)}^q}
            \\
            &\le \E{L^q \norm{X_n - X}^{\alpha q}}
            = L^q \norm{f(X_n) - f(X)}_{L^p}^p \to 0.
            \qedhere
        \end{align*}
    \end{solution}
\end{exercise}

\begin{exercise}[Counterexample \(L^p\) continuous mapping \Coffeecup]
    Let \(f(x) = x^\gamma\) with \(\gamma >1\), let \(U\sim \uniform(0,1)\) and define
    \[
        X_n = n^{\frac1q} \ind{[0, \frac1n)}(U).
    \]
    Show for \(p\in [\frac{q}\gamma, q)\) we have that \(X_n \overset{L^p}\to 0\)
    but not \(f(X_n) \overset{L^p}\to 0\).
    \begin{solution}
        We have
        \[
            \norm{X_n-0}_{L^p}^p = \norm{X_n}_{L^p}^p = n^{\frac{p}{q}} \prob{U < \tfrac1n} = n^{\frac{p}{q}-1} \to 0
        \]
        for \(p < q\). However, we have
        \[
            \norm{f(X_n)- f(0)}_{L^p}^p = \norm{f(X_n)}_{L^p}^p = n^{\frac{\gamma p}{q}} \prob{U < \tfrac1n} = n^{\frac{\gamma p}{q}-1} \to \infty
        \]
        for \(p > \frac{q}\gamma\).    
    \end{solution}
\end{exercise}

\begin{exercise}[Basic application of continuous mapping theorem \Coffeecup]
    Let \((X_n)_{n\in \nat}\) and \(X\) be real-valued random variables
    with \(X_n \overset{d}\to X\). Show that
    \begin{enumerate}[label=(\alph*), noitemsep]
        \item \(X_n^2 \overset{d}\to X^2\),
        \begin{solution}
            This is a direct application of the continuous mapping theorem with the function \(f(x) = x^2\) that is continuous everywhere.            
        \end{solution}

        \item \(\sgn(X_n) \overset{d}\to \sgn(X)\) if \(\prob{X=0} = 0\).

        Construct a counterexample for the case \(\prob{X=0} > 0\).
        \begin{solution}
            For the first part we again apply the continuous mapping theorem with the function \(f(x) = \sgn(x)\) that is continuous at all points except zero. Since \(\prob{X=0} = 0\) we have that \(f\) is almost surely continuous in \(X\) and thus the continuous mapping theorem applies.

            For the counterexample we can take any deterministic sequence \(x_n\to 0\)
            such that \(\sgn(x_n)\) does not converge to \(\sgn(0)\). This does depend a bit on the definition
            of \(\sgn(0)\) but exists in all cases. Then \(X_n = x_n\) has the same problem.
            Specifically, for the Lipschitz function \(g(x) = \min\set{\abs{x},
            1}\) we have
            \[
                \E{g(X_n)} = x_n \not\to 0 = \E{g(0)}
            \]
            and thus \(X_n \overset{d}{\not\to} 0\) by \ref{it: b.lip}.
        \end{solution}
    \end{enumerate}
\end{exercise}

\begin{exercise}[Slutsky's theorem \Coffeecup]
    Let \((X_n)_{n\in \nat}\) and \(X\) be random variables in \(\real\)
    such that \(X_n \overset{d}\to X\). And let \((Y_n)_{n\in \nat}\) be random
    variables in \(\real\) such that \(Y_n \overset{d}\to c\) for some
    constant \(c\in \real\). Prove
    \begin{enumerate}[label=(\alph*)]
        \item \(X_n + Y_n \overset{d}\to X + c\),
        \item \(X_n Y_n \overset{d}\to X c\),
        \item if \(c \neq 0\) then \(X_n/Y_n \overset{d}\to X/c\).
    \end{enumerate}
    \begin{solution}
        Since \(X_n \overset{d}\to X\) and \(Y_n \overset{d}\to c\) implies
        \((X_n, Y_n) \overset{d}\to (X,c)\) by Theorem \ref{thm: joint
        convergence}~\ref{it: joint conv law slutsky}, we have all the claims by
        the continuous mapping theorem, since the sum and product are continuous
        and \(\frac1x\) is continuous in \(c \neq 0\).
    \end{solution}
\end{exercise}

\begin{exercise}[Find the mistake \Coffeecup]
    Let \((X_n)_{n\in \nat}\) and \(X\) be random variables in \(\real^\dims\).
    Which step in the following chain of implications is \textbf{wrong}?
    \[
        X_n \overset{d}\to X
        \overset{?}\implies X_n - X \overset{d}\to 0
        \overset{?}\implies X_n - X \overset{p}\to 0
        \overset{?}\implies X_n \overset{p}\to X.
    \]
    Justify all other steps and find a counterexample for the wrong step.
    \begin{solution}
        The first step is wrong. Consider \(X_n = -X\) with
        \(\prob{X=-1}=\prob{X=1}=\frac12\). Then
        as in  Exercise \ref{ex: convergence in distribution but not in
        probability} \(X_n \to X\), but \(X_n - X = -2X\) does not converge to
        zero in distribution.

        All the other steps are correct. The second step is the special case
        that convergence in distribution to a constant implies convergence in
        probability to that constant (Theorem \ref{thm: convergence
        implications}~\ref{it: d to p}). The third step follows from 
        \[
            \prob{\norm{X_n - X} > \epsilon}
            =\prob[\big]{\norm{(X_n - X) - 0} > \epsilon}
            \to 0.
            \qedhere
        \]
    \end{solution}
\end{exercise}

\begin{exercise}[Incompatible metric \Coffeecup]
    \label{ex: incompatible metric}
    Let \((x_n)_{n\in \nat}\) and \((y_n)_{n\in \nat}\) be sequences in
    \(\real\). Assume that \(x_n \to x\) and \(y_n \to y\) in the usual
    Euclidean metric.
    Show that this does not necessarily imply \((x_n, y_n) \to (x,y)\)
    in \((\real^2, \metric_\infty)\) with
    \[
        \metric_\infty((x,y), (x',y'))
        = \begin{cases}
            0 & x=y\\
            1 & x\neq y.
        \end{cases}
    \]
    \begin{solution}
        With the counterexample \(x_n = \frac1n\) and \(y_n = \frac1n\) we have
        that \(x_n \to 0\) and \(y_n \to 0\) but \((x_n, y_n)\) does not
        converge to \((0,0)\) in \((\real^2, \metric_\infty)\) since
        \[
            \metric_\infty((x_n, y_n), (0,0)) = 1
        \] for all \(n\).
    \end{solution}
\end{exercise}